%% ===================================================================
\section{Leading logarithms: BFKL and JIMWLK}\label{sec:LL}
%% ===================================================================

\subsection{Where the logarithms come from}\label{sec:LLorigin}

The unobserved (or virtual) gluon has $\Lambda<p^+<\vee-k^+$. As explained in Sec.~\ref{sec:rows}, a large logarithm arises only where its momentum is strongly ordered with respect to the measured $k^+$. There are two such regions (Fig.~\ref{fig:rapidity}):
\begin{itemize}
\item \textbf{Region T}, $\Lambda<p^+\ll k^+$. The extra gluon is softer than the measured one. It is emitted by \emph{all} harder colour charges: the valence charges $\rho$ \emph{and} the measured gluon. Its logarithm is $\delta Y_T=\log(k^+/\Lambda)$.
\item \textbf{Region P}, $k^+\ll p^+<\vee$. The extra gluon is harder. Now the measured gluon is the softer one and is emitted by the valence charges \emph{and} by the extra gluon. The logarithm is $\delta Y_P=\log(\vee/k^+)$.
\end{itemize}
At LL accuracy the full interval $\delta Y=\log(\vee/\Lambda)$ is therefore shared, $\delta Y_T+\delta Y_P=\delta Y$. It cannot multiply both evolution kernels, because plus momentum is conserved at every splitting and a gluon can only be emitted by a parent with larger $p^+$. The BFKL ``real'' term, in which the measured gluon is emitted by the extra gluon, needs $p^+>k^+$. Conversely, the JIMWLK terms in which the soft gluon is emitted by the measured gluon need $p^+<k^+$.

\begin{figure}[t]
\centering
\begin{tikzpicture}[>=Stealth,scale=1.0]
\draw[->,thick] (0,0) -- (12.4,0) node[right] {$\log p^+$};
\foreach \x/\lab in {0.6/\Lambda, 6.0/k^+, 11.4/\vee}{
  \draw[thick] (\x,0.15) -- (\x,-0.15);
  \node[below] at (\x,-0.2) {$\lab$};}
\fill[blue!12] (0.6,0.3) rectangle (5.0,1.1);
\fill[red!12] (7.0,0.3) rectangle (11.4,1.1);
\node at (2.8,0.7) {\small region T: target (JIMWLK)};
\node at (9.2,0.7) {\small region P: projectile (BFKL)};
\draw[decorate,decoration={brace,amplitude=6pt,mirror}] (0.6,-0.75) -- (6.0,-0.75) node[midway,below=7pt] {$\delta Y_T=\log\frac{k^+}{\Lambda}$};
\draw[decorate,decoration={brace,amplitude=6pt,mirror}] (6.0,-0.75) -- (11.4,-0.75) node[midway,below=7pt] {$\delta Y_P=\log\frac{\vee}{k^+}$};
\node[below] at (0.6,-1.7) {\small target};
\node[below] at (11.4,-1.7) {\small projectile};
\node[below] at (6.0,-1.7) {\small measured gluon};
\end{tikzpicture}
\caption{The rapidity interval of the soft window. The unobserved gluon gives a large logarithm only when it is much softer (region T) or much harder (region P) than the measured gluon. At mid rapidity both $\delta Y_T$ and $\delta Y_P$ are large.}
\label{fig:rapidity}
\end{figure}

\paragraph*{The leading-log wave function.} Denote by
\begin{equation}
j^a_q(s)\equiv\sum_i a^{b\dagger}_i(q^+,s)\,(T^a)_{bc}\,a^c_i(q^+,s)
\label{eq:jdef}
\end{equation}
the colour charge density of the gluons with plus momentum $q^+$. It obeys the same algebra as $\rho$, $[j^a(x),j^b(y)]=if^{abc}j^c(x)\delta^{(2)}(x-y)$, and rotates in the same way, $S^\dagger j^a(x)S=U^{ab}(x)j^b(x)$ (using $U^TT^aU=U^{ab}T^b$). The emission of gluons with momentum $q^+$ by a colour source $J$ is generated by
\begin{equation}
\Omega_q[J]=\exp\Big\{i\varepsilon_q\sum_{e,j}\int_{z,s}\KK^j(s-z)\,J^e(s)\,\big(a^e_j(q^+,z)+a^{e\dagger}_j(q^+,z)\big)\Big\},\qquad \varepsilon_q^2\ \leftrightarrow\ \frac{\alpha_s}{\pi^2}\,\frac{dq^+}{q^+},
\label{eq:OmegaLL}
\end{equation}
with the modes normalized per bin of $\log q^+$. The normalization of $\varepsilon_q$ is that of Eq.~\eqref{eq:A1X}: to first order $\Omega_q[\rho]|0\rangle$ reproduces Eq.~\eqref{eq:LOwf}, and each emission costs $\frac{\alpha_s}{\pi^2}\frac{dq^+}{q^+}$, as noted after Eq.~\eqref{eq:LO}. At LL the full $\Omega$ factorizes into strongly ordered emissions, the harder gluon first:
\begin{equation}
\text{region T:}\quad \Omega=\Omega_p[\rho+j_k]\,\Omega_k[\rho],\qquad\qquad \text{region P:}\quad \Omega=\Omega_k[\rho+j_p]\,\Omega_p[\rho].
\label{eq:OmegaLLfact}
\end{equation}
The coefficients of Sec.~\ref{sec:omega} reproduce Eq.~\eqref{eq:OmegaLLfact} in the ordered limits. For example, Eq.~\eqref{eq:B2soft} is the order-$\varepsilon_k\varepsilon_p$ two-gluon term of $\Omega_p[\rho+j_k]\Omega_k[\rho]$. Its $\KK^j(w-z)$ term is the $j_k$ part of the source, and its $\frac12\KK^j(x-z)$ term is the commutator part of the product of the two $\rho$-emissions. Likewise $\mf C_1$ for $p^+\ll k^+$ (its $\delta_{jm}\delta_{ik}$ term) is the soft gluon emitted by the measured gluon after the target. The unitarity relation~\eqref{eq:unitB1} holds for Eq.~\eqref{eq:OmegaLLfact}.

\subsection{Region T: the JIMWLK Hamiltonian}\label{sec:T}

\paragraph*{Step 1: the cross section as an expectation value of the LO operator.} Insert $\Omega=\Omega_p[J]\,\Omega_k$, with $J\equiv\rho+j_k$, into Eq.~\eqref{eq:xsec}. With $N\equiv a^{a\dagger}_i(k^+,w)a^a_i(k^+,w')$ and $SS^\dagger=1$,
\begin{equation}
\langle0|\Omega^\dagger S^\dagger\Omega N\Omega^\dagger S\Omega|0\rangle=\langle\Psi|\,W^\dagger\,\bar M\,W\,|\Psi\rangle,\qquad W\equiv\bar\Omega_p^\dagger\Omega_p,\quad|\Psi\rangle\equiv\Omega_k|0\rangle,\quad\bar M\equiv S^\dagger\Omega_kN\Omega_k^\dagger S.
\label{eq:step1}
\end{equation}
Without the soft gluon, $W=1$, and $\langle\Psi|\bar M|\Psi\rangle$ is the LO cross section: it is the operator
\begin{equation}
\Ord(w,w')\equiv\int_{x,x'}\KK(x-w)\cdot\KK(x'-w')\,\rho^b(x)\Big[\big(\Ud(x)-\Ud(w)\big)\big(U(x')-U(w')\big)\Big]^{bb'}\rho^{b'}(x')
\label{eq:Oop}
\end{equation}
of Eq.~\eqref{eq:LOop}.

\paragraph*{Step 2: emission before and after the target.} Write $\Omega_p=e^{i\varepsilon G}$ with $G=\sum_{h,j}\int_zA^h_j(z)\,\alpha^h_j(z)$, where $\alpha\equiv a(p^+)+a^\dagger(p^+)$ and
\begin{equation}
A^h_j(z)\equiv\int_s\KK^j(s-z)\,J^h(s).
\end{equation}
Using $S^\dagger J^e(s)S=U^{eg}(s)J^g(s)$ and $S^\dagger\alpha^e(z)S=U^{eh}(z)\alpha^h(z)$, the rotated generator is $\bar G=\sum\int_zB^h_j(z)\,\alpha^h_j(z)$ with
\begin{equation}
B^h_j(z)\equiv\int_s\KK^j(s-z)\,\big[U^T(z)U(s)J(s)\big]^h=\int_s\KK^j(s-z)\,U^{eh}(z)U^{eg}(s)J^g(s).
\end{equation}
$A$ is the emission of the soft gluon before the target and $B$ its emission after the target. In both, the soft colour index $h$ refers to the frame before the target.

\paragraph*{Step 3: real and virtual terms.} To order $\varepsilon^2$,
\begin{equation}
W=e^{-i\varepsilon\bar G}e^{i\varepsilon G}=1+i\varepsilon(G-\bar G)-\frac{\varepsilon^2}{2}\big(G^2+\bar G^2\big)+\varepsilon^2\bar GG+\Ord(\varepsilon^3).
\end{equation}
$|\Psi\rangle$ contains no soft gluon, and $\langle0_p|\alpha^h_j(z)\alpha^{h'}_{j'}(z')|0_p\rangle=\delta^{hh'}\delta_{jj'}\delta^{(2)}(z-z')$. Write $X\cdot Y\equiv\sum_{h,j}\int_zX^h_j(z)Y^h_j(z)$ and
\begin{equation}
\gamma\equiv A-B .
\end{equation}
The order-$\varepsilon^2$ terms of Eq.~\eqref{eq:step1} are
\begin{align}
\text{REAL}={}&\varepsilon^2\,\langle\Psi|\,\gamma\cdot\bar M\,\gamma\,|\Psi\rangle,\label{eq:real}\\
\text{VIRTUAL}={}&\varepsilon^2\,\Big\langle\Psi\Big|\Big[-\tfrac12(A\cdot A+B\cdot B)+A\cdot B\Big]\bar M+\bar M\Big[-\tfrac12(A\cdot A+B\cdot B)+B\cdot A\Big]\Big|\Psi\Big\rangle.\label{eq:virt}
\end{align}
REAL has the soft gluon in the final state; it is the region-T part of Groups~II and~III. VIRTUAL has one gluon in the final state; it is the region-T part of Group~I. The sources $s$ in $A$ and $B$ run over the valence charges ($J=\rho$) and over the measured gluon ($J=j_k$).

\paragraph*{Step 4: real $+$ virtual is a double commutator.} Since
\begin{equation}
-\tfrac12[\gamma\cdot,[\gamma\cdot,\bar M]]=\gamma\cdot\bar M\gamma-\tfrac12\{A\cdot A+B\cdot B,\bar M\}+\tfrac12(A\cdot B+B\cdot A)\bar M+\tfrac12\bar M(A\cdot B+B\cdot A),
\end{equation}
we find
\begin{equation}
\text{REAL}+\text{VIRTUAL}=-\frac{\varepsilon^2}{2}\,\Big\langle\Psi\Big|\,[\gamma\cdot,[\gamma\cdot,\bar M]]+[[B\cdot,A\cdot],\bar M]\,\Big|\Psi\Big\rangle .
\label{eq:dc}
\end{equation}
The weights $(1,-2,1)$ of $\gamma\gamma\bar M$, $\gamma\bar M\gamma$ and $\bar M\gamma\gamma$ are those of the virtual terms, the real terms and the conjugate virtual terms.

\paragraph*{Step 5: colour charges act as derivatives on the Wilson lines.} Define the right and left rotation operators on functionals of $U$~\cite{JalilianMarian:1997gr,Kovner:2005nq},
\begin{equation}
J_R^a(u)\,U(x)=\delta^{(2)}(x-u)\,U(x)T^a,\qquad J_L^a(u)\equiv U^{ab}(u)J_R^b(u),\qquad J_L^a(u)\,U(x)=\delta^{(2)}(x-u)\,T^aU(x).
\label{eq:JRJL}
\end{equation}
\textbf{Lemma.} For every operator $X$ that does not depend on $U$,
\begin{equation}
J_R^a(u)\big(S^\dagger XS\big)=\big[S^\dagger XS,\,J^a(u)\big],\qquad J_L^a(u)\big(S^\dagger XS\big)=\big[S^\dagger XS,\,\bar J^a(u)\big],\qquad \bar J^a(u)\equiv U^{ab}(u)J^b(u).
\label{eq:dual}
\end{equation}
\emph{Proof.} $S^\dagger XS$ is built from the rotated fields $S^\dagger\rho^b(x)S=U^{bc}(x)\rho^c(x)$ and $S^\dagger a^b(x)S=U^{bc}(x)a^c(x)$. Both sides of Eq.~\eqref{eq:dual} obey the Leibniz rule, so it suffices to check them on these fields. For $\rho$, $J_R^a(u)\big(U^{bc}(x)\rho^c(x)\big)=\delta^{(2)}(x-u)\big(U(x)T^a\big)^{bc}\rho^c(x)$, while $\big[U^{bc}(x)\rho^c(x),\rho^a(u)\big]=U^{bc}(x)\,if^{cad}\rho^d(x)\,\delta^{(2)}(x-u)=\delta^{(2)}(x-u)U^{bc}(x)(T^a)_{cd}\rho^d(x)$. The two are equal. The gluon fields work in the same way with $j$. The $J_L$ form follows by multiplying with $U^{ab}(u)$. $\square$

Write $Y^g(s)\equiv J_R^g(s)\bar M=[\bar M,J^g(s)]$ and $Z^m(s)\equiv J_L^m(s)\bar M=[\bar M,\bar J^m(s)]$. Then $Z^m(s)=S^\dagger[X,J^m(s)]S$ has the form required by the lemma, while $Y^g(s)=U^{mg}(s)S^\dagger[X,J^m(s)]S$ contains an explicit $U(s)$, so that
\begin{equation}
[Y^g(s'),J^{g'}(s)]=J_R^{g'}(s)\,Y^g(s')-\delta^{(2)}(s-s')\,(T^{g'})_{kg}\,Y^k(s').
\label{eq:extra}
\end{equation}

\paragraph*{Step 6: the double commutators.} With $A^h=\int_s\KK(s-z)J^h(s)$, $B^h=\int_s\KK(s-z)U^{mh}(z)\bar J^m(s)$ and the shorthand $\iint\equiv\int_z\int_{s,s'}\KK(s-z)\cdot\KK(s'-z)$:
\begin{align}
[A,[A,\bar M]]&=\iint J_R^h(s)J_R^h(s')\,\bar M,\\
[B,[B,\bar M]]&=\iint J_L^m(s)J_L^m(s')\,\bar M,\\
[A,[B,\bar M]]&=\iint U^{mh}(z)\,J_R^h(s)J_L^m(s')\,\bar M,\\
[B,[A,\bar M]]&=\iint U^{mh}(z)\,J_L^m(s)J_R^h(s')\,\bar M+\int_z\int_s\KK(s-z)^2\,\Tr\big[\Ud(z)U(s)T^k\big]J_R^k(s)\,\bar M,\\
[[B,A],\bar M]&=\int_z\int_s\KK(s-z)^2\,\Tr\big[\Ud(z)U(s)T^k\big]J_R^k(s)\,\bar M.
\end{align}
In the first line the extra term of Eq.~\eqref{eq:extra} is proportional to $(T^h)_{kh}=0$. The trace term in the fourth line comes from Eq.~\eqref{eq:extra} with $U^{mh}(z)U^{mg}(s)(T^g)_{kh}=-\Tr[\Ud(z)U(s)T^k]$. In the last line we used $[B^h,A^h]=\int_s\KK(s-z)^2U^{mh}(z)U^{mg}(s)\,if^{ghk}J^k(s)$. In the combination $[\gamma,[\gamma,\bar M]]+[[B,A],\bar M]$ the two trace terms cancel, and
\begin{align}
\text{REAL}+\text{VIRTUAL}={}&-\frac{\varepsilon^2}{2}\int_z\int_{s,s'}K(s,s',z)\Big\{J_R^h(s)J_R^h(s')+J_L^m(s)J_L^m(s')\notag\\
&\hspace{9em}-U^{mh}(z)\big[J_R^h(s)J_L^m(s')+J_L^m(s)J_R^h(s')\big]\Big\}\langle\Psi|\bar M|\Psi\rangle .
\label{eq:HLR}
\end{align}
$|\Psi\rangle$ does not depend on $U$, so the derivatives act on the LO operator~\eqref{eq:Oop}. The points $s,s'$ run over the positions of its Wilson lines: the valence points $x,x'$ and the gluon points $w,w'$. A derivative at $w$ or $w'$ is the soft gluon emitted by the measured gluon.

\paragraph*{Step 7: identification with JIMWLK.} The JIMWLK Hamiltonian~\cite{JalilianMarian:1997jx,JalilianMarian:1997gr,JalilianMarian:1998cb,Kovner:2000pt,Weigert:2000gi,Iancu:2000hn,Ferreiro:2001qy} is
\begin{equation}
H_{\rm JIMWLK}=\frac{\alpha_s}{2\pi^2}\int_{t,u,v}K(u,v,t)\,\big[\Ud(u)-\Ud(t)\big]_{ce}J_R^c(u)\,\big[U(v)-U(t)\big]_{ed}J_R^d(v),
\label{eq:HJIMWLK}
\end{equation}
with every derivative acting on everything to its right. Using $[\Ud(u)]_{ce}J_R^c(u)=J_L^e(u)$, the integrand is $\big(J_L^e(u)-U^{ec}(t)J_R^c(u)\big)\big(J_L^e(v)-U^{ed}(t)J_R^d(v)\big)$, which expands to the bracket of Eq.~\eqref{eq:HLR}. With $\varepsilon^2=\frac{\alpha_s}{\pi^2}\delta Y_T$ we obtain
\begin{equation}
\boxed{\ \frac{d^3N}{d^3k}\bigg|_{\rm T}=-\,\delta Y_T\;H_{\rm JIMWLK}\,\frac{d^3N_{\rm LO}}{d^3k},\qquad \delta Y_T=\log\frac{k^+}{\Lambda}.\ }
\label{eq:Tresult}
\end{equation}
This is the target evolution $\partial_Y=-H_{\rm JIMWLK}$ of the LO spectrum over the interval below the measured gluon. When all $J_R$ in Eq.~\eqref{eq:HJIMWLK} are moved to the right of the Wilson lines, a term linear in $J_R$ appears; it is included automatically in Eq.~\eqref{eq:HLR}. The explicit action on Eq.~\eqref{eq:LO} is given in App.~\ref{app:explicit}.

\subsection{Region P: the BFKL equation}\label{sec:P}

Now $\Omega=\Omega_k[J']\,\Omega_p$ with $J'\equiv\rho+j_p$. The number operator $N$ involves only gluons with momentum $k^+$, so $[\Omega_p,N]=0$ and
\begin{equation}
\Omega N\Omega^\dagger=\Omega_k[J']\,N\,\Omega_k^\dagger[J'].
\end{equation}
The cross section is therefore the LO calculation with the source $\rho$ replaced by $\rho+j_p$, evaluated in the evolved projectile $|\Phi\rangle\equiv\Omega_p|0\rangle$:
\begin{equation}
\langle\Ord\rangle\ \longrightarrow\ \big\langle\Phi\big|\,\Ord[\rho+j_p]\,\big|\Phi\big\rangle,\qquad|\Phi\rangle=\Big(1+i\varepsilon G-\tfrac{\varepsilon^2}{2}G^2+\dots\Big)|0\rangle,\qquad G=\int_{u,z}\KK^j(u-z)\rho^e(u)\,\alpha^e_j(z).
\end{equation}
Write $\Ord[J']=\int_{x,y}\KK(x-w)\cdot\KK(y-w')\,J'^b(x)\,\mathbb M^{bb'}(x,y)\,J'^{b'}(y)$ with $\mathbb M(x,y)\equiv(\Ud(x)-\Ud(w))(U(y)-U(w'))$. There are three types of order-$\varepsilon^2$ terms.

\paragraph*{(i) The $\rho\rho$ part.} This is the evolution of the charge correlator itself, with real and virtual parts,
\begin{align}
&\varepsilon^2\big\langle G\rho\mathbb M\rho G-\tfrac12\{G^2,\rho\mathbb M\rho\}\big\rangle\notag\\
&\qquad=-\frac{\varepsilon^2}{2}\int_z\int_{u,u'}K(u,u',z)\,\big[\rho^h(u),\big[\rho^h(u'),\rho^b(x)\rho^{b'}(y)\big]\big]\,\mathbb M^{bb'}.
\end{align}
Using $(if^{hbc})(if^{hcd})=N_c\delta_{bd}$ and $(if^{hbc})(if^{hb'd})=(T^h)_{bc}(T^h)_{b'd}$,
\begin{align}
\text{(i)}={}&-\frac{\varepsilon^2}{2}\int_z\Big\{N_c\big[K(x,x,z)+K(y,y,z)\big]\rho^b(x)\rho^{b'}(y)\notag\\
&\hspace{5em}+2K(x,y,z)\,(T^h)_{bc}(T^h)_{b'd}\,\rho^c(x)\rho^d(y)\Big\}\mathbb M^{bb'}.
\end{align}
For a colour-diagonal correlator, $(T^h)_{bc}(T^h)_{b'c}=-N_c\delta_{bb'}$, and this becomes
\begin{equation}
-\frac{\varepsilon^2}{2}\,N_c\,K_D(x,y,z)\,\frac{\langle\rho^f(x)\rho^f(y)\rangle}{N_c^2-1}\,\Tr\,\mathbb M,
\end{equation}
with the dipole kernel
\begin{equation}
K_D(x,y,z)\equiv K(x,x,z)-2K(x,y,z)+K(y,y,z)=\frac{(x-y)^2}{(x-z)^2(y-z)^2}.
\label{eq:KD}
\end{equation}

\paragraph*{(ii) The $j_pj_p$ part: the measured gluon is emitted by the harder gluon.} With $j^{b'}(y)a^{e\dagger}(z)|0\rangle=\delta^{(2)}(y-z)(T^{b'})_{ge}a^{g\dagger}(z)|0\rangle$,
\begin{equation}
\text{(ii)}=\varepsilon^2\langle G\,j\,\mathbb M\,j\,G\rangle=\varepsilon^2\int_z\int_{u,u'}K(u',u,z)\,\rho^{e'}(u')\rho^e(u)\,(T^bT^{b'})_{e'e}\,\KK(z-w)\cdot\KK(z-w')\,\mathbb M^{bb'}(z,z).
\end{equation}
After the colour average, $\Tr(T^bT^{b'})=N_c\delta^{bb'}$.

\paragraph*{(iii) The $\rho\,j_p$ cross terms.}
\begin{equation}
\varepsilon^2\langle G\,\rho^b(x)\mathbb M^{bb'}j^{b'}(y)\,G\rangle=\varepsilon^2\int_z\int_{u,u'}K(u',u,z)\,\delta^{(2)}(y-z)\,\rho^{e'}(u')\rho^b(x)\rho^e(u)\,(T^{b'})_{e'e}\,\mathbb M^{bb'}(x,z).
\end{equation}
$K(u',u,z)$ is symmetric and $(T^{b'})_{e'e}$ is antisymmetric under $(e',u')\leftrightarrow(e,u)$, so the fully symmetric three-charge part \emph{vanishes exactly}. Only commutators survive:
\begin{equation}
\rho^{e'}(u')\rho^b(x)\rho^e(u)(T^{b'})_{e'e}\to\tfrac12N_c\,\delta^{(2)}(u-u')\,\rho^{b'}(u)\rho^b(x)+\delta^{(2)}(x-u)\,f^{beg}f^{b'e'e}\,\rho^{e'}(u')\rho^g(x).
\end{equation}
After the colour average this gives $\varepsilon^2N_c\KK(x-w)\cdot\KK(z-w')\big[\frac12K(y,y,z)-K(x,y,z)\big]$ times the trace with $\mathbb M(x,z)$, and the mirror term with $x\leftrightarrow y$, $w\leftrightarrow w'$.

\paragraph*{Result.} With $\varepsilon^2=\frac{\alpha_s}{\pi^2}\delta Y_P$, the sum (i)$+$(ii)$+$(iii) is the BFKL evolution of the projectile charge correlator in the LO spectrum:
\begin{equation}
\boxed{\ \frac{d^3N}{d^3k}\bigg|_{\rm P}=\delta Y_P\,\big[{\rm BFKL}\otimes{\rm LO}\big],\qquad \delta Y_P=\log\frac{\vee}{k^+},\ }
\label{eq:Presult}
\end{equation}
where ${\rm BFKL}\otimes{\rm LO}$ is Eq.~\eqref{eq:LOBFKL} of App.~\ref{app:explicit}. It is obtained by applying the coordinate-space BFKL equation for $\langle\rho^f(x)\rho^f(y)\rangle$, Eq.~\eqref{eq:vBFKL}, to the LO spectrum~\eqref{eq:LO}. Term (i) gives the line with $K_D$, (ii) the line with $\KK(z-w)\cdot\KK(z-w')$, and (iii) the two lines with $2K(x,y,z)-K(y,y,z)$ and $2K(x,y,z)-K(x,x,z)$. Term (i) contains Group~I (the $G^2$ normalization) and part of Groups~II and~III; (ii) and (iii) belong to Group~II.

\subsection{Cancellation of three- and four-charge terms}\label{sec:cancel}

\paragraph*{Region T.} The total is $-\delta Y_TH_{\rm JIMWLK}\Ord$. $H_{\rm JIMWLK}$ contains only $J_R$, $J_L$ and $U$, which act on Wilson lines and never change the number of charges. $\Ord$ has two charges, so the result has exactly two. Every three- and four-charge term of the real rows is cancelled by the corresponding virtual term. The mechanism is visible in Eq.~\eqref{eq:dc}.
\begin{itemize}
\item \emph{Four charges.} For the valence part of $\gamma$, treat the charges as commuting numbers. Then $\gamma\gamma\bar M-2\gamma\bar M\gamma+\bar M\gamma\gamma=0$.
\item \emph{Three charges.} Take the soft-gluon factor of a spectator charge, i.e.\ a charge not contained in $\bar M$. It commutes with $\bar M$ and with the other $\gamma$, so its double commutator vanishes. The real term $\varepsilon^2\gamma_{\rm spec}\{\bar M,\gamma_{\rm other}\}$ is cancelled exactly by the virtual term $-\varepsilon^2\gamma_{\rm spec}\{\bar M,\gamma_{\rm other}\}$. Only commutators with the charges of $\bar M$ survive, and they leave two charges.
\end{itemize}
\paragraph*{Region P.} Term (i) is a double commutator of two charges, which gives two charges, and (ii) has two. In (iii) the symmetric three-charge part vanishes by the antisymmetry shown above.

\subsection{Row-by-row leading logarithms}\label{sec:LLrows}

The derivation above uses only the factorized wave function~\eqref{eq:OmegaLLfact}. It can be compared row by row with Sec.~\ref{sec:xsec}. In each region, every coefficient reduces to a product of Weizs\"acker--Williams factors, as in Eqs.~\eqref{eq:B2soft} and~\eqref{eq:B2hard}. Inserting these limits into Eqs.~\eqref{eq:CS1}--\eqref{eq:CS3} gives the LL content of every row. The two-gluon amplitudes $\mathfrak B$, $\mathfrak X_{1,2}$, $\mathfrak C$ and the one-gluon corrections $\mathcal V_X$ at LL are listed in App.~\ref{app:LLrows}. Two points deserve mention.
\begin{itemize}
\item In Group~I the unitarity rewriting of $\bar{\mf A}^{(3)\dagger}_1$ distributes one physical term over Rows G1-II to G1-V. At LL their sum is a single simple term, listed as G1-II(b) in App.~\ref{app:LLrows}. It contains the commutator $[\bar{\mf N}_2,\bar{\mf A}^{(1)}]$ of Eq.~\eqref{eq:VII}.
\item Row G1-I$'$, the soft gluon emitted and reabsorbed after the target by a measured gluon that was produced before the target, is needed. Without it the parts of $H_{\rm JIMWLK}\Ord$ with the kernels $K(w,w,z)$, $K(w',w',z)$, $K(x',w,z)$ and $K(x,w',z)$ are not reproduced.
\end{itemize}
We fully symmetrized all products of charges (Weyl ordering plus commutator terms) and compared the rows term by term. In region T, all 12 classes of four-charge terms and all 63 classes of three-charge terms cancel among the rows. Of the two-charge classes, 137 cancel and 40 survive. In region P, 12 four-charge and 43 three-charge classes cancel; of the two-charge classes 105 cancel and 32 survive. Grouped by kernel, the survivors are exactly Eq.~\eqref{eq:HLR} in region T (8 kernels) and the terms (i)--(iii) in region P (8 kernels). The checks are described in App.~\ref{app:checks}.

\subsection{Result at mid rapidity}\label{sec:LLresult}

Adding the two regions,
\begin{equation}
\boxed{\ \frac{d^3N}{d^3k}\bigg|_{\rm LL}=\log\frac{\vee}{k^+}\,\big[{\rm BFKL}\otimes{\rm LO}\big]+\log\frac{k^+}{\Lambda}\,\big[{\rm JIMWLK}\otimes{\rm LO}\big],\ }
\label{eq:LLmain}
\end{equation}
with ${\rm JIMWLK}\otimes{\rm LO}\equiv-H_{\rm JIMWLK}\,\frac{d^3N_{\rm LO}}{d^3k}$ and BFKL$\otimes$LO given by Eq.~\eqref{eq:LOBFKL}.
This is the expected rapidity factorization of single inclusive production~\cite{Kovchegov:2001sc,Baier:2005dv,Kovner:2006ge,Kovner:2006wr,Gelis:2008rw}, here obtained directly from the $\Ord(g^4)$ cross section. At mid rapidity both logarithms are large and both evolutions must be resummed: the projectile charge correlator over $\log(\vee/k^+)$ and the target over $\log(k^+/\Lambda)$.

Using the first identity of Eq.~\eqref{eq:logid}, Eq.~\eqref{eq:LLmain} can be rewritten as
\begin{equation}
\frac{d^3N}{d^3k}\bigg|_{\rm LL}=\log\frac{\vee}{\Lambda}\,\big[{\rm JIMWLK}\otimes{\rm LO}\big]+\log\frac{\vee}{k^+}\,\Big(\big[{\rm BFKL}\otimes{\rm LO}\big]-\big[{\rm JIMWLK}\otimes{\rm LO}\big]\Big).
\end{equation}
So the coefficient of the total interval $\log(\vee/\Lambda)$ is JIMWLK$\otimes$LO alone. BFKL$\otimes$LO is the coefficient of $\log\vee$ at fixed $\Lambda$ and $k^+$. It would be incorrect to compare the coefficient of $\log(\vee/\Lambda)$ with the sum BFKL$\otimes$LO$+$JIMWLK$\otimes$LO.
